\section*{Ensayo por Alexa Reynoso Ascencio}
\addcontentsline{toc}{subsubsection}{Ensayo por Alexa Reynoso Ascencio.}

Dado el crecimiento significativo de la tecnología en los últimos años, el cual ha impulsado nuevos avances y formas de entretenimiento, existe un aspecto sumamente importante que solemos ignorar y que resulta fundamental en la actualidad: la información se ha consolidado como un elemento clave del poder dentro de la economía digital. Considero que este tema es de suma relevancia debido a que, como lo menciona la autora Anahiby Becerril y como muchos de nosotros hemos experimentado, las empresas recopilan y comercializan nuestros datos personales diariamente, sin nuestro consentimiento plenamente informado, aprovechándose del desconocimiento general que existe sobre el verdadero valor de nuestra información.

La autora destaca cómo el Big Data y el Internet de las Cosas (IoT) se han convertido en los pilares fundamentales de la economía digital. Por un lado, el Big Data hace referencia a conjuntos de datos masivos y complejos cuyo procesamiento no se puede realizar mediante herramientas tradicionales de gestión de bases de datos; por otro lado, el IoT actúa como una estructura global que interconecta dispositivos para transmitir información en tiempo real. Esto se relaciona directamente con la materia de Fundamentos de Bases de Datos, ya que, aunque comúnmente los datos se organizan en esquemas estructurados y definidos antes de su almacenamiento, el enorme volumen y la velocidad con la que se captura la información en estos casos exige diseños de almacenamiento mucho más amplios, dinámicos y eficientes. Debido a que al final del día, las empresas dependen de esta infraestructura para categorizar, procesar y extraer valor de la información personal, construyendo así mejores servicios e interfaces más adictivas para los usuarios.

A pesar de que no toda la población cuenta aún con acceso a internet, es un hecho que con el paso del tiempo esta cifra seguirá aumentando y, por ende, la cantidad de información generada crecerá exponencialmente. Es por esto que la autora enfatiza la necesidad de ser conscientes de cuánta información compartimos diariamente, ya sea de manera voluntaria o involuntaria, a través de nuestros dispositivos electrónicos. Mediante el análisis de nuestros datos personales, las empresas diseñan estrategias publicitarias altamente personalizadas y obtienen beneficios basados en el contenido generado por los propios usuarios. Hoy en día, la mayoría de nosotros hemos experimentado el impacto de estas estrategias; por ejemplo, al notar cómo los algoritmos de las redes sociales se adaptan a nuestras preferencias para mantenernos más tiempo enganchados a una aplicación, o al recibir sugerencias de compras en línea basadas en nuestras búsquedas recientes. Todo esto recalca la importancia de tomar conciencia sobre la economía de los datos personales y exigir una mayor transparencia en el uso y tratamiento de los mismos.

Por otro lado, la autora aborda cómo, aunque no exista un método único para asignar un costo exacto a la información personal, existe una clara diferencia entre la valoración del mercado y la percepción individual sobre la privacidad. Mientras que las grandes empresas generan ganancias multimillonarias monetizando nuestros datos, cuando a un usuario se le llega a recompensar por ellos, se le ofrece una miseria en comparación con el valor real que representan. 

Lo que la autora busca dar a entender es que nuestros datos son, en realidad, el pago implícito que realizamos para acceder a servicios que aparentemente son "gratuitos". Aunque a veces intentamos ser conscientes de lo que compartimos, la mayoría de las ocasiones no le damos la atención necesaria; un caso claro en el que la mayoría podemos identificarnos es al aceptar los términos y condiciones de una plataforma sin leerlos o comprender realmente su alcance y lo que significa. Además, últimamente se ha intensificado la necesidad de extraer datos cada vez más íntimos: un ejemplo son los relojes o anillos inteligentes que rastrean parámetros de salud, horas de sueño, ejercicio y localización, o las aplicaciones para el seguimiento del ciclo menstrual, donde se registra información altamente sensible sobre la salud reproductiva e íntima de las mujeres. Aunque estos sistemas nos facilitan llevar el control, muchas empresas terminan utilizando esta información sin un consentimiento verdaderamente informado para realizar sus propios estudios y negocios. 

Por lo tanto, desde mi perspectiva como desarrolladora, las empresas tenemos la responsabilidad de establecer con absoluta claridad cómo se almacenará, procesará y protegerá la información que se nos está proporcionando, evitando problemas éticos o legales y brindando confianza a los usuarios. Pero también como usuarios queda en nosotros ser más consciente de que información aceptaremos compartir y cuál no.

En conclusión, tal como lo señala la autora Anahiby Becerril, la economía digital ha transformado de manera radical el uso y la importancia que las empresas otorgan a nuestra información personal. Esto reafirma la postura de que debemos tomar mayor conciencia sobre los datos que proporcionamos diariamente y las repercusiones directas que su uso puede tener sobre nuestra privacidad. En lo personal, aunque este tema es algo que se me había mencionado en otras ocasiones, el análisis del artículo me permitió profundizar y comprender con claridad por qué ocurre esto con nuestros datos. Finalmente, todo esto nos crea retos futuros cruciales, como si a futuro el precio para acceder a estos servicios continuará siendo nuestra información personal, o si veremos cambios hacia un modelo de recopilación de información más transparente, ético y respetuoso con la información de las personas. 
